\subsubsection*{IMERG: promedio exacto de cuenca y validación (V07B)}
Actualización del 4 de octubre de 2026. IMERG Final Run V07B (GPM\_3IMERGM), variable Grid/precipitation, malla de 0,1$^\circ$. Se verificaron 300 archivos mensuales de enero de 1998 a diciembre de 2022. El promedio se calcula sobre el polígono CAMELS de La Vieja con 42 celdas intersectadas, 31 de borde, ponderadas por el área geodésica WGS84 de su intersección. Cobertura espacial 100.000000\%. Área geodésica del polígono 2780.18 km²; el área CAMELS de 2797,19 km² usada para R permanece sin cambios. Las tasas originales están en mm/h y se multiplican por las horas de cada mes, incluidos los bisiestos. No se rellenan celdas ni se renormalizan pesos cuando faltan valores. Hay 300/300 meses válidos y 285 meses completos comunes con CHIRPS y Q. Enero de 1998: 0.10202291 mm/h $\times$ 744 h = 75.905043 mm/mes; se comprobó también la suma independiente celda por celda. El promedio Giovanni de caja y sus resultados previos se conservan como antecedentes y no se mezclan con esta nueva serie. Los bordes se segmentizan a máximo 0,001$^\circ$ para las áreas geodésicas. Se reconstruyen límites contiguos de la malla nominal 0,1$^\circ$ para evitar los solapes del redondeo float32 de los límites originales; ambos límites quedan registrados. No se modifica ni remuestrea la precipitación. La suma de intersecciones coincide con el área del polígono dentro de una tolerancia relativa de una parte por millón. La malla es de unos 11 km por lado en esta región: aporta representación espacial de la cuenca, pero no resuelve toda la variabilidad orográfica interna.

Acceso: GES DISC mediante el script de descarga conservado y sesión Earthdata. Reproducción: \path{scripts/22_imerg_poligono_estadisticas.py}; huellas, pesos y metadatos en \path{documentos/imerg_poligono}.

DOI: \url{https://doi.org/10.5067/GPM/IMERG/3B-MONTH/07}.
\begin{figure}[H]\centering\includegraphics[width=\linewidth]{figuras/imerg_poligono_series.pdf}\caption{Fuentes identificadas y alineadas, 1998--2022; los vacíos se conservan.}\end{figure}
